\scanpage{53}
\mspage{14}
On suppose
\[
(23)\left\{\begin{aligned}
\text{a)}\quad &X_{\Delta_1}\text{ est un \uline{sous-espace} \inserted{top} \uline{fermé} de }X_{\Delta_0}\times X_{\Delta_0}\\
&\text{(avec la top. induite)}\\
\text{b)}\quad &\forall(i,j)\in I\times I,\quad\Gamma_{ij}\overset{\mathrm{def}}{=}X_{\Delta_1}\cap(X_i\times X_j)\quad\text{\deleted{\unsure{en}} est}\\
&\text{\inserted{ou bien vide, ou bien}}\\
&\text{le graphe d'une immersion fermée}\\
&\deleted{\text{de}}\quad X_i\deleted{\text{\unsure{ds}}}\longrightarrow X_j
\quad\text{(donc }X_{\Delta_1}\xrightarrow{\sigma}X_{\Delta_0}
\end{aligned}\right.
\]
est un morphisme de rev. \inserted{étale} constant, $\}$
\deleted{\unsure{et alors|étale}} et
$X_{\Delta_1}\overset{b}{\unsure{\longrightarrow|-}}X_{\Delta_0}$ une imm. locale.)

Alors la relation $(i\leq j)\overset{\mathrm{def}}{\Longleftrightarrow}\Gamma_{ij}\ne\varnothing$ est une
relation d'ordre sur $I$, \deleted{\unsure{les $X_j$|les $X_i$}} et les $X_i$
(avec les $X_i\unsure{-|\longrightarrow}X_j$ de b) forment un
système inductif d'espaces topologiques,
indexé par $I$, avec morphismes de
transition des imm. fermées, et
la donnée de l'espace \inserted{top} ordonné $X_{\Delta_0}$,
\deleted{et} avec la partition envisagée de $X_{\Delta_0}$,
revient à la donnée d'un tel syst.
inductif, indexé par un ens ordonné $I$.
\deleted{Supposons}\inserted{Supposons} de plus pour nous mettre
à l'aise que $I$ est fini, et que les
$\operatorname{Inf}$ des deux éléments y existe --- alors
(sauf erreur) $X=\underset{i\in I}{\unsure{\underset{\longrightarrow}{\mathrm{Lim}}|\underline{\mathrm{Lim}}}}\ X_i\simeq
X_{\Delta_0}/(X_{\Delta_1}\rightrightarrows)$
et les $X_i\hookrightarrow X$ sont des imm fermées
etc. ...
\margin{52}
